\uuid{xDvr}
\titre{Image réciproque}
\auteur{bodin ; hanani ; mzari}
\organisation{exo7}
\licence{CC BY-NC-SA 4.0}
\datecreate{2020-01-20}
\niveau{L1}
\module{Algèbre}
\chapitre{Logique, ensemble, raisonnement}
\sousChapitre{Ensemble}
\difficulte{3}
\theme{Ensembles, applications, image réciproque}
\exo7qcmid{lille-1.ensembles.023}

\contenu{
\begin{question}
\qclassification{100.02,101.01,102.01,102.02}
\qtags{motcle=image réciproque}
Soient $E$ et $F$ deux ensembles non vides et $f$ une application de $E$ dans $F$. Soit $A$ un sous-ensemble de $E$. Quelles sont les bonnes réponses ?
\begin{answers}
    \bad{$A=f^{-1}(f(A))$}
    \good{$A\subset f^{-1}(f(A))$}
    \bad{$f^{-1}(f(A))\subset A$}
    \bad{$f^{-1}(f(A))=E\setminus A$}
\end{answers}
\begin{explanations}
Pour tout $x\in A$, on a $f(x)\in f(A)$, donc $x\in f^{-1}(f(A))$.
\end{explanations}
\end{question}
}
